% Source: https://artofproblemsolving.com/community/c3h2621680_basic_number_theory
% Stable identity: aops-topic-2621680-q01
Consider the $p-1$ numbers $a,2a,\dots,(p-1)a$.
\begin{itemize}
\item None is divisible by $p$: $p$ is prime and divides neither $k$ (for $1\le k\le p-1$) nor $a$.
\item No two are congruent mod $p$: $ia\equiv ja$ would give $p\mid(i-j)a$, so $p\mid i-j$ and $i=j$.
\end{itemize}
So modulo $p$ they are $1,2,\dots,p-1$ in some order. Multiplying them all,
\[a^{p-1}(p-1)!\equiv(p-1)!\pmod p.\]
Since $p\nmid(p-1)!$, we can cancel it and get $a^{p-1}\equiv1\pmod p$. $\blacksquare$

The argument used only $p\nmid a$, never the size or sign of $a$. (One can also reduce $a$ mod $p$ to a remainder $r$ with $1\le r\le p-1$, and $a^{p-1}\equiv r^{p-1}$.)

The condition $p\nmid a$ is needed: if $p\mid a$ then $a^{p-1}\equiv0\not\equiv1\pmod p$. For example $7^6\equiv0\pmod7$.
